% Part: normal-modal-logic
% Chapter: tableaux
% Section: simple-S5

\documentclass[../../../include/open-logic-section]{subfiles}

\begin{document}

\olfileid{nml}{tab}{s5}

\olsection{\Log{S5} لئی سادے \usetoken{P}{tableau}}

\Log{S5} ہر پاسے رسائی والے نمونیاں دی جماعت لئی صحت مند تے مکمل اے،
یعنی اوہ نمونے جیہناں وچ ہر جہان ہر جہان توں رسائی وچ اے۔ ایسے
نمونیاں وچ رسائی دا تعلق فرق نہیں پاؤندا۔ «کوئی جہان~$w$ اے جتھے
$\mSat{M}{!A}[w]$» اُدوں تے صرف اُدوں سچ اے جدوں اجیہا کوئی $w$،
~$u$ توں رسائی وچ اے۔ سو \Log{S5} وچ نمونیاں دی تعریف لئی جہاناں
دا اک سیٹ تے قدر بندی~$V$ ای کافی نیں۔ ایس توں لگدا اے کہ اسیں
!!{tableau} دے قاعدے وی سوکھے کر سکدے آں۔ عام صورت وچ اگیتے مثبت
صحیح عدداں دیاں لڑیاں ہوندے نیں، تاں جو پتا رہے کہ کہڑے اگیتے اوہناں
جہاناں دے ناں نیں جیہڑے دوجے جہاناں توں رسائی وچ نیں۔ $\sigma.n$
ایسے جہان دا ناں اے جیہڑا~$\sigma$ توں رسائی وچ اے۔ پر \Log{S5} وچ
ہر جہان ہر جہان توں رسائی وچ اے، ایس لئی اوہ حساب رکھن دی لوڑ نہیں۔
اگیتیاں لئی مثبت صحیح عدد ورت سکدے آں۔ سوکھے قاعدے
\olref{tab:rules-S5} وچ دتے نیں۔

\begin{table}
  \begin{center}
    \def\arraystretch{3}\def\fCenter{}
    \begin{tabular}{|c|c|}
    \hline
    \iftag{prvBox}{
      \AxiomC{\sFmla{\True}{\Box !A}[n]}
      \RightLabel{\TRule{\True}{\Box}}
      \UnaryInfC{\sFmla{\True}{!A}[m]}
      \DisplayProof
      &
      \AxiomC{\sFmla{\False}{\Box !A}[n]}
      \RightLabel{\TRule{\False}{\Box}}
      \UnaryInfC{\sFmla{\False}{!A}[m]}
      \DisplayProof\\
      $m$ ورتیا ہویا اے & $m$ نواں اے
      \\[1ex]
      \hline}{}
    \iftag{prvDiamond}{
      \AxiomC{\sFmla{\True}{\Diamond !A}[n]}
      \RightLabel{\TRule{\True}{\Diamond}}
      \UnaryInfC{\sFmla{\True}{!A}[m]}
      \DisplayProof
      &
      \AxiomC{\sFmla{\False}{\Diamond !A}[n]}
      \RightLabel{\TRule{\False}{\Diamond}}
      \UnaryInfC{\sFmla{\False}{!A}[m]}
      \DisplayProof\\
      $m$ نواں اے & $m$ ورتیا ہویا اے
      \\[1ex]
      \hline}{}
    \end{tabular}
  \end{center}
  \caption{\Log{S5} لئی سوکھے قاعدے۔}
  \ollabel{tab:rules-S5}
\end{table}

\begin{ex}
  اسیں اک سادہ بند تابلو دیندے آں جیہڑا وکھاؤندا اے کہ $\Log{S5} \Proves
  \Ax{5}$، یعنی $\Diamond!A \lif \Box\Diamond!A$.
  \begin{oltableau}
    [\pFmla{\False}{\Diamond\formula{A} \lif \Box\Diamond \formula{A}}{1},
      just = \TAss
      [\pFmla{\True}{\Diamond \formula{A}}{1}, just = {\TRule{\False}{\lif}[1]}
        [\pFmla{\False}{\Box\Diamond \formula{A}}{1},
          just = {\TRule{\False}{\lif}[1]}
          [\pFmla{\False}{\Diamond \formula{A}}{2},
            just = {\TRule{\False}{\Box}[3]}
            [\pFmla{\True}{\formula{A}}{3},
              just = {\TRule{\True}{\Diamond}[2]}
                [\pFmla{\False}{\formula{A}}{3},
                  just = {\TRule{\False}{\Diamond}[4]}, close]
            ]
          ]
        ]
      ]
    ]
  \end{oltableau}
\end{ex}

\end{document}
